\documentclass[10pt,a4paper]{amsart}
\usepackage[margin=19mm]{geometry}
\usepackage{amsmath,amssymb,mathtools}
\usepackage[scr=rsfs,cal=euler]{mathalfa}
\usepackage{xfrac}
\usepackage[unicode,hidelinks,bookmarks=false]{hyperref}
\newcommand{\ie}{i.e.\ }
\newcommand{\cf}{cf.\ }
\newcommand{\resp}{respectively\ }
\newcommand{\Subsection}{\subsection}
\providecommand{\nobreakdash}{\nobreak\hbox{-}}
\setlength{\emergencystretch}{2em}
\multlinegap0pt
\usepackage{amssymb}
\newtheorem{thm}{Theorem}
\DeclareMathOperator{\Fr}{Fr}
\DeclareMathOperator{\Span}{Span}
\newcommand{\V}{\mathcal{V}}
\newcommand{\field}{\mathbb{F}}
\newcommand{\invring}{S_m^{G_m}}
\title{Invariants of Finite Orthogonal Groups of Plus Type in Odd Characteristic}
\author{H. E. A. Campbell, R. James Shank, David L. Wehlau}
\date{Source: arXiv:2407.01152v3 (CC BY 4.0)}
\begin{document}
\maketitle

\noindent\emph{Source and licence.} Invariants of Finite Orthogonal Groups of Plus Type in Odd Characteristic. Version arXiv:2407.01152v3 (CC BY 4.0), distributed by arXiv under the Creative Commons Attribution 4.0 International licence, http://creativecommons.org/licenses/by/4.0/, as recorded in the arXiv licence metadata for this exact version and shown on the abstract page https://arxiv.org/abs/2407.01152v3. Full licence notice: \texttt{licenses/arxiv-2407.01152-CC-BY-4.0.txt}.\par\smallskip

\noindent\textit{Editorial scope.} Counted unit: the article's thm orth\_var (source0.tex lines 287--291). The article's complete original proof of it is delivered with it (source0.tex lines 292--350, one \texttt{proof} environment of 59 lines). No mathematics is written, repaired or re-derived: the statement and the proof are the article's own lines, in the article's own notation, and no retained line is substituted or re-flowed.  No further source-local context is needed: every cross-reference of every retained line resolves inside the two delivered spans.  Nothing was omitted from any retained span, so the package carries no omission record.  No retained line contains a citation, so the wrapper carries no bibliography and no bibliography entry.  No retained line contains a figure environment, an image-inclusion command or drawing code, so no image is delivered and the package declares no asset dependency.  Editorial formatting only, in addition: an AMS \texttt{amsart} preamble that restates the article's own declarations and macros used by the retained lines, namely the environments \texttt{thm} on separate counters, together with the standard packages those lines need.  The retained environments print in this document as Theorem 1 in order of appearance, whereas the article prints the same items with the numbering of its own full document.  The complete native source of the article as distributed in the arXiv e-print for this exact version is bundled unchanged as \texttt{sources/161619/source0.tex}.
\begin{thm}\label{orth_var} For $s\geq m$,
%updated  m-1 to m here
$$\V(\xi_1,\ldots,\xi_s)=\bigcup_{g\in G_m}
g\cdot\Span_{\overline{\field}_q}\{e_1,\ldots,e_m\}.$$
\end{thm}

\begin{proof} 
Observe that the variety $\V(\xi_1,\ldots,\xi_s)$ is closed under scalar multiplication 
since it is defined by homogeneous polynomials.
Let $W_m$ denote 
$\Span_{\overline{\field}_q}\{e_1,\ldots,e_m\}$.
Since $\xi_{i}(v)=0$ for $v\in W_m$ and $\xi_{i}\in \invring$, 
$$\bigcup_{g\in G_m}g\cdot W_m\subseteq \V(\xi_1,\ldots,\xi_s).$$ 
Therefore, it is sufficient to prove that for 
$v\in \V( \xi_1,\ldots,\xi_{m})$, there exists $g\in G_m$ such that
$gv\in W_m$.

The proof is by induction on $m$. For $m=1$, we have $x_1(v)y_1(v)=0$. If $x_1(v)\not=0$, then take $g$ to be the transposition which interchanges $x_1$ and $y_1$.

For $m>1$, define $\bar{v}:=v-y_m(v)e_m-x_m(v)f_m$, let
$\bar{\xi_{i}}$ denote the restriction of $\xi_{i}$ to 
$\Span_{\overline{\field}_q}\{e_1,\ldots,e_{m-1},f_{m-1},\ldots f_1\}$  and
identify $G_{m-1}$ with the point-wise stabiliser of  
${\rm Span}_{\field_q}\{e_m,f_m\}$ in $G_m$.

If $x_m(v)=0$ then $\bar{\xi_{i}}(\bar{v})=0$ for $i=1,\ldots,m$.
By induction, there is an element $g\in G_{m-1}< G_m$
with $g\bar{v}\in W_{m-1}$. Observe that 
$gv\in W_m$.


For $u\in \overline{V}$ we use $\Fr^i(u)$ to denote the $i^{th}$ iteration of the Frobenius map on $u$.
Recall that $\beta(u,u)=2\xi_1(u)$ and note that, for $i\geq2$,
$\beta(u,\Fr^{i-1}(u))=\xi_{i}(u)$.

If $x_m(v)\not=0$ then we scale $v$ so that $x_m(v)=1$ and define $w:=v-\Fr(v)$.
Observe that
\begin{eqnarray*}
\xi_1(w)&=&
\beta(v-\Fr(v),v-\Fr(v))/2
=\left(\beta(v,v)+\beta(\Fr(v),\Fr(v))-2\beta(v,\Fr(v))\right)/2\\
   &=&\xi_1(v)+\xi_1(v)^q-\xi_2(v)=0
\end {eqnarray*}
and, for $i\geq2$,
\begin{eqnarray*}
\xi_{i}(w)&=&\beta(v-\Fr(v),\Fr^{i-1}(v-\Fr(v)))=\beta(v-\Fr(v),\Fr^{i-1}(v)-\Fr^{i}(v))\\
&=&\beta(v,\Fr^{i-1}(v))-\beta(\Fr(v),\Fr^{i-1}(v))-\beta(v,\Fr^{i}(v))+\beta(\Fr(v),\Fr^{i}(v))\\
   &=&\xi_{i}(v)-\xi_{i-1}(v)^q-\xi_{i+1}(v)+\xi_{i}(v)^q=0.
\end {eqnarray*}
Therefore $\xi_{i}(w)=0$ for $i=1,\ldots,m-1$.
Since $x_m(w)=0$,
we have $\bar{\xi_{i}}(\bar{w})=0$ for $i=1,\ldots, m-1$ and so by induction
there is an element $g\in G_{m-1}<G_m$
with $g\bar{w}\in W_{m-1}$ and $gw\in W_m$.
Hence $(x_j-x_j^q)(gv)=x_j(gw)=0$ for $j=1,\ldots,m$ and $x_j(gv)\in\field_q$.
Since $g$ stabilises ${\rm Span}_{\field_q}\{e_m,f_m\}$, we have $x_m(gv)=1$.
For convenience, define $c_j=x_j(gv)\in\field_q$ for $j=1,\ldots,m-1$ and
let $h$ denote the linear transformation
given by $y_jh=y_j$ for $j=1,\ldots,m-1$, 
and $y_mh=y_m+\sum_{j=1}^{m-1} c_jy_j$, and $x_jh=x_j-c_jx_m$ for $j=1,\ldots,m-1$
and $x_mh=x_m$. Observe that $h\in G_m$
and, since  $\xi_1(gv)=0$, we have $y_m(hgv)=y_m(gv)+\sum_{j=1}^{m-1} c_jy_j(gv)=0$.
Let $\sigma $ denote the transposition which exchanges $x_m$ and $y_m$. We then have
$x_m(\sigma h g v)=0$ and we can apply the induction argument as above.
\end{proof}

\end{document}
